We calculate zero-shot perplexity on the Penn Tree Bank (PTB) \cite{marcus1994penn} dataset measured in \cite{radford2019language}. We omit the 4 Wikipedia-related tasks in that work because they are entirely contained in our training data, and we also omit the one-billion word benchmark due to a high fraction of the dataset being contained in our training set.  PTB escapes these issues due to predating the modern internet. Our largest model sets a new SOTA on PTB by a substantial margin of 15 points, achieving a perplexity of 20.50. Note that since PTB is a traditional language modeling dataset it does not have a clear separation of examples to define one-shot or few-shot evaluation around, so we measure only zero-shot.


\begin{table}
    
    \centering
        \begin{tabular}{ll}
        \toprule
        Setting & PTB  \\ 
        \midrule
        SOTA (Zero-Shot) & 35.8\textsuperscript{\textit{a}} \\
        GPT-3 Zero-Shot & \textbf{20.5}\\
        \bottomrule
        \end{tabular}
        
    \caption{\textbf{Zero-shot results on PTB language modeling dataset.} Many other common language modeling datasets are omitted because they are derived from Wikipedia or other sources which are included in GPT-3's training data. \textsuperscript{\textit{a}}\cite{radford2019language}}
    \label{table:language}
\end{table}